\chapter[KESIMPULAN DAN SARAN]{Kesimpulan dan Saran}
\label{chap:kesimpulan-saran}

Bab ini merangkum hasil perancangan dan eksperimen OMNI-RAG pada \emph{User Manual} dan
dokumen hukum. Kesimpulan berikut menjawab tiga rumusan masalah berdasarkan \emph{Recall}@30.
Setiap konfigurasi dihitung secara terpisah tanpa penggabungan hasil antarkonfigurasi.

\section{Kesimpulan}
\label{sec:kesimpulan}

\begin{enumerate}
  \item Arsitektur \emph{plugin} mendukung empat implementasi
  \emph{chunker}, yaitu \emph{generic structure-aware}, \emph{generic recursive},
  \emph{generic sliding-window}, dan \emph{legal structure-aware}, tanpa mengubah kode inti.
  \emph{Sliding-window} digunakan sebagai \emph{baseline}. Pada \emph{User Manual}, konfigurasi
  \emph{Structure-aware hybrid} memperoleh \emph{Recall}@30 sebesar 0,740, sedangkan \emph{baseline}
  \emph{Sliding-window} memperoleh 0,685, dengan selisih +0,055. Pada dokumen hukum,
  \emph{Legal-structure dense} memperoleh \emph{Recall}@30 sebesar 0,6854, sedangkan \emph{baseline}
  \emph{Sliding} memperoleh 0,6345, dengan selisih +0,0509. Dengan demikian, pemeliharaan struktur dokumen meningkatkan cakupan
  \emph{context reference} pada kedua \emph{use case}.

  \item Arsitektur \emph{plugin} mendukung tiga strategi
  \emph{indexer}, yaitu \emph{sparse}, \emph{dense}, dan \emph{hybrid}, tanpa mengubah orkestrator.
  \emph{Structure-aware sparse} digunakan sebagai \emph{baseline}. Pada \emph{User Manual},
  \emph{Structure-aware dense} berbasis BAAI/\emph{BGE-M3} memperoleh \emph{Recall}@30 sebesar
  0,785, sedangkan \emph{baseline sparse} memperoleh 0,420. Pada dokumen hukum, \emph{Legal-structure
  dense} memperoleh 0,6854, sedangkan \emph{baseline sparse} memperoleh 0,6381. Dengan demikian,
  pencarian \emph{dense} menjadi strategi terbaik untuk tujuan utama penelitian, yaitu menemukan
  sebanyak mungkin \emph{context reference}. \(K = 30\) digunakan sebagai cutoff bersama karena
  kurva \emph{recall} mulai memasuki bagian yang lebih datar pada titik tersebut.

  \item Arsitektur \emph{plugin} mendukung dua strategi
  \emph{knowledge graph}, yaitu \emph{legal graph} dan \emph{LLM graph}, dengan artefak kanonik
  yang sama dan tanpa perubahan pada orkestrator. Eksperimen aktif menggunakan \emph{legal graph}
  pada dokumen hukum dan menelusuri \emph{neighborhood} global sampai kedalaman dua. Pada
  \(K = 30\), \emph{knowledge-graph recall} mencapai 1,0000 untuk \emph{Sliding}, 0,9975 untuk
  \emph{Recursive} dan \emph{hybrid}, 0,9672 untuk \emph{sparse}, serta 0,9556 untuk
  \emph{Legal-structure dense}. Nilai ini menjawab jangkauan dokumen terhubung dan dibaca
  terpisah dari \emph{context-reference recall}. Dengan demikian, seluruh rumusan masalah
  terjawab melalui kontrak \emph{plugin}, hasil \emph{recall} per komponen, dan traversal graf
  yang dapat ditelusuri.
\end{enumerate}

\section{Saran}
\label{sec:saran}

\noindent Pengembangan berikutnya dapat dilakukan melalui beberapa arah berikut.

\begin{enumerate}
  \item \textbf{Memperluas \emph{dataset} dan korpus.} Uji strategi pada lebih banyak manual,
  peraturan, bahasa, dan pertanyaan terverifikasi. Pertahankan pemetaan \emph{context reference} serta
  pembagian strata agar perbandingan antarkorpus tetap adil.

  \item \textbf{Memisahkan tahap \emph{retrieval}, \emph{reranking}, dan \emph{generation}.}
  Pertahankan \(K = 30\) sebagai cutoff kandidat saat menguji \emph{reranker}. Untuk
  \emph{generation}, laporkan \(K = 10\), bahasa, model penilai, dan jumlah pertanyaan agar
  hasil dapat direproduksi.

  \item \textbf{Memperkuat eksperimen komponen.} Uji ukuran segmen, \emph{overlap}, model
  \emph{embedding}, strategi \emph{indexer}, dan pengaturan konteks induk secara terpisah sebelum
  menyimpulkan pengaruh gabungannya.

  \item \textbf{Memperdalam pengukuran \emph{knowledge graph}.} Pertahankan pemisahan antara
  \emph{context-reference recall} dan \emph{neighborhood recall}. Simpan jenis relasi, arah relasi,
  kedalaman traversal, dan dokumen yang membuka setiap tetangga agar kontribusi graf dapat
  dianalisis per pertanyaan.

  \item \textbf{Meningkatkan kesiapan operasional \emph{plugin}.} Sediakan paket instalasi yang
  memuat \emph{manifest}, versi kontrak, dependensi, pemeriksaan kompatibilitas, dan pengujian
  untuk setiap \emph{plugin}. Validasi ini mencegah artefak, indeks, dan graf yang tidak kompatibel
  ketika sumber atau strategi pemrosesan diganti.
\end{enumerate}
